\section{Reakcje podporowe w układach płaskich}

\subsection*{Warunek równowagi}

Reakcje wyznaczamy z warunku równowagi statycznej:

\begin{equation}
    \begin{cases}
        \Vec{F}_0 = \Vec{0} \\
        \Vec{M}_0 = \Vec{0}
    \end{cases}
\end{equation}
który w płaskim układzie sił przyjmuje postać:

\begin{equation}
    \begin{cases}
        F_{0,x} = 0 \\
        F_{0,y} = 0 \\
        M_{0,z} = 0
    \end{cases}
\end{equation}

\subsection*{Układ dwóch sił i momentu skupionego}

Jeżeli reakcje stanowi układ dwóch wektorów sił i jegnego momentu skupionego:
\begin{itemize}
    \item Równanie momentów rozpisujemy względem punktu $A$, w którym przecinają się osie działania wektorów sił - z tego równania wyznaczamy wartość momentu skupionego,
    \item Wartości sił uzyskujemy z równań zerowania się rzutów sił na kierunkach $x$ i $y$.
\end{itemize}

\begin{equation}
    \begin{cases}
        \sum M_z^A = 0 \\
        \sum F_x = 0 \\
        \sum F_y = 0
    \end{cases}
\end{equation}

Jeżeli reakcje stanowi układ trzech wektorów sił:
\begin{itemize}
    \item Równanie momentów rozpisujemy względem punktów $A$ i $B$, w których przecinają się osie działania wektorów sił - z tych równania wyznaczamy dwie z trzech niewiadomych sił,
    \item Trzecią siłę uzyskujemy z równania zerowania się rzutów sił na kierunku $x$ lub $y$,
    \item Ostatnie równanie (suma rzutów sił na kierunku $x$ lub $y$) traktujemy jako równanie sprawdzające.
\end{itemize}

\begin{equation}
    \begin{cases}
        \sum M_z^A = 0 \\
        \sum M_z^B = 0 \\
        \sum F_x = 0 \\
        \sum F_y = 0
    \end{cases}
\end{equation}